\section{模型的改进与推广}

\subsection{模型的改进}
可将当前"存储感知邻域回退搜索"升级为"任务—储能联合 LNS-MILP"：把任务起点离散为候选集后，与储能变量一起进入小规模 MILP，以四格解为热启动做局部联合优化，从而给出更紧的可行成本；远期可为 $I$ 提供有效下界（如 LP 松弛或拉格朗日松弛），把点估计升格为认证区间。

预测侧可引入概率预测（分位数回归/双层）与滚动时域，把"对未来负荷的确定性追求"替换为对边界情景的鲁棒保护，缓解高噪声序列的信息量不足。

\subsection{模型的推广}
本框架可推广到含多时间尺度（日前—日内—实时）的算力市场出清、含绿色电力证书/碳交易市场的联合调度，以及含任务优先级与故障迁移的在线调度；当把单位 GPU 功率与 PUE 换成实测动态曲线后，可直接用于真实数据中心群的运行辅助决策。
